\section{Metodologia di sviluppo}
In questa sezione si descriverà l'organizzazione del team di sviluppo.

\subsection{Scrum}
La metodologia di sviluppo applicata su questo progetto è di tipo \textit{Agile}, in particolare
si è adottata la variante \textit{Scrum}.

Il processo di sviluppo Scrum è un processo \textit{iterativo} e \textit{incrementale}, in cui a ogni
iterazione si aggiungono nuove funzionalità al sistema o si raffinano delle funzionalità
pre-esistenti.

La prima attività all'interno del processo di sviluppo Scrum è quella di redarre il
\textit{Product Backlog}, ovvero una lista di funzionalità, dette \textit{items}.

Le iterazioni del processo, chiamate \textit{sprint}, saranno svolte settimanalmente. Ogni sprint
prevede le seguenti attività:
\begin{itemize}
    \item \textit{Sprint Planning}: riunione iniziale in cui si selezionano gli item da implementare
    durante lo sprint e si scompongono in sotto-funzionalità, di cui il team di sviluppo
    fornisce stime sulla loro complessità e alcuni design di dettaglio. Al termine dello
    Sprint Planning si produce uno \textit{Sprint Backlog} che assegna a ogni membro del team
    dei \textit{task} da eseguire.
    \item \textit{Daily Scrum}: una breve riunione giornaliera, in cui il team di sviluppo si aggiorna
    sul progresso del progetto, riadattando lo Sprint Backlog.
\end{itemize}

Al termine dello sprint si prevedono tre ulteriori attività:
\begin{itemize}
    \item \textit{Product Backlog Refinement}: una riunione in cui il team di sviluppo si accorda sulle
    migliorie che devono essere apportate al Product Backlog.
    \item \textit{Sprint Review}: si valuta il progresso del progetto ottenuto al termine dello sprint,
    verificando che sia un \textit{PSPI (Potentially Shippable Product Increment)}.
    \item \textit{Sprint Retrospective}: si valuta il processo di sviluppo adottato, discutendo possibili
    cambiamenti che potrebbero aumentare l'efficacia del team.
\end{itemize}

L'organizzazione del personale all'interno del processo Scrum prevede tre ruoli:
\begin{itemize}
    \item \textit{Product Owner}: si occupa di redarre il \textit{Product Backlog} e di verificare l'adeguatezza del
    sistema realizzato. Il ruolo è stato assunto da Jahrim Gabriele Cesario;
    \item \textit{Scrum Master}: agisce da esperto del processo Scrum e da mediatore tra gli altri due ruoli.
    Il ruolo è stato assunto da Mirko Felice;
    \item \textit{Development Team}: si occupa di progettare soluzioni adeguate ai task definiti dal Product
    Owner, stimando tempi di realizzazione e proponendo modifiche sul sistema al Product Owner.
    Il team di sviluppo sarà composto da:
    \begin{itemize}
        \item Remedi Tommaso
        \item Biagioni Giacomo
        \item Morsucci Federico
        \item Morbidelli Cristian

    \end{itemize}
\end{itemize}

Le riunioni saranno svolte periodicamente su \textit{Microsoft Teams}, mentre l'organizzazione
delle funzionalità del progetto sarà tenuta su \href{https://trello.com/b/2mgr8e1j/pps}{\textit{Trello}}.

\subsection{Test-Driven Development}
Durante lo sviluppo del sistema è stato scelto di applicare il più possibile il \textit{Test-Driven Development (TDD)},
il cui scopo è quello di anticipare il prima possibile la fase di testing per minimizzare i costi di manutenzione e
il rischio di fallimento del progetto.

Il processo seguito durante il TDD è un processo iterativo chiamato \textit{Red-Green-Refactor (RGR)}, che prevede a ogni
iterazione le seguenti fasi:
\begin{enumerate}
    \item \textit{Red}: scrivere un test che fallisca per una certa funzionalità da implementare
    \item \textit{Green}: scrivere il codice di produzione che soddisfi il test definito precedentemente
    \item \textit{Refactor}: ristrutturare sia il codice di testing che quello di produzione
\end{enumerate}

A supporto di questo processo, sono stati adottati i seguenti strumenti:
\begin{itemize}
    \item \textit{ScalaTest}: framework per la definizione di unit test per scala
    \item \textit{SCoverage}: strumento per valutare la qualità dei test come percentuale di codice di produzione analizzato
\end{itemize}

Poiché si è deciso di utilizzare il testing solo per quanto riguarda il modello, è stato scelto di utilizzare lo stile
di testing \textit{FlatSpec}, in quanto questo stile è un buon compromesso rispetto alla semplicità di \textit{FunSuite} e
alla organizzazione dei test di \textit{FunSpec}. Inoltre, permette di realizzare dei test in modo dichiarativo e quindi più
vicini all'utente finale.

\subsection{Quality Assurance}
Per il controllo della qualità del sistema sono stati adottati i seguenti strumenti:
\begin{itemize}
    \item \textit{Scala Formatter}: controlla lo stile del codice scala
    \item \textit{Wart Remover}: individua possibili difetti nel codice scala
    \item \textit{Ktlint}: controlla lo stile del codice kotlin
    \item \textit{Detekt}: individua possibili difetti nel codice kotlin
\end{itemize}

\subsection{Build Automation}
Per automatizzare i processi di compilazione, testing e release del codice sviluppato si è deciso di utilizzare \textit{Gradle}
come strumento di \textit{Build Automation}.

È stato preferito Gradle al posto di Sbt, perché è una tecnologia più matura e gli sviluppatori hanno più esperienza
con tale strumento.

\subsection{Continuous Integration}
Per fare in modo che il codice rimanga integro e corretto durante lo sviluppo, è stato utilizzato un workflow che
attraverso \textit{GitHub Actions} permette di eseguire i test del progetto su diverse configurazioni di sistemi operativi,
a ogni aggiornamento del progetto.

\subsection{Continuous Delivery}
Analogamente, un altro workflow permette di automatizzare la release su \textit{GitHub Releases} e \textit{Maven Central}.

\subsection{Automated Evolution}
Per mantenere aggiornate le dipendenze del progetto in maniera automatica è stato utilizzato \textit{Renovate Bot}, che osserva
le dipendenze del progetto e in caso di necessità apre delle \textit{pull request} sul repository GitHub per aggiornarle.

\subsection{Versioning}
Per il versionamento del sistema, è stato adottato lo standard \textit{Semantic Versioning}, il quale prevede che una versione
sia caratterizzata da:
\begin{itemize}
    \item \textit{Major}: numero identificativo da aggiornare a ogni cambiamento non-retrocompatibile
    \item \textit{Minor}: numero identificativo da aggiornare a ogni aggiunta o modifica retrocompatibile di una funzionalità del
    sistema
    \item \textit{Patch}: numero identificativo da aggiornare a ogni correzione dei difetti del sistema
\end{itemize}

Adottando lo standard \textit{Conventional Commits} durante lo sviluppo tramite DVCS, è stato possibile
automatizzare l'assegnamento della versione al sistema sulla base dei commit effettuati. Per fare ciò, è stato sfruttato
\textit{Semantic Release Bot} per GitHub.